\section{Language Independent Storage Strategies}

%\subsection{The rstrategies Framework}
%\todo{Framework description}

%\subsection{Storage Strategies for Objects}
%\todo{RSqueak Implementation}
%\todo{which strategies are most important? why?}

%\subsubsection{Single Value Strategies}
%\todo{Nil}

%\subsubsection{Single Type Strategies}
%\todo{Int, Float}

%\subsubsection{Mixed Type Strategies}
%\todo{NilAndInt, NilAndFloat}

%\subsection{Supporting Exotic Objects}
%\todo{ContextParts, Classes, MethodDicts}


In this chapter we present the design of our library for storage strategies called rstrategies.
First we introduce the principles this library design is based on, then we describe different parts of the library and how a VM developer can utilize the library to create specific collection implementations.

In order to be reusable in the context of different programming languages, rstrategies must expose its functionality in a very flexible way.
We use the concept of mixins \cite{brachaMixins, mixins} to provide the required level of flexibility while allowing the VM to extend the provided functionality when needed.

In the context of RPython, any regular class can be used as a mixin.
The programmer must simply include a \lstinline{import_from_mixin(MixinClass)} clause into the body of the instantiating class.
All attributes and methods will then be copied over from \lstinline{MixinClass} to the target class, including all inherited methods.
A mixin is not intended to be used on its own.
It can require other mixins and fields or methods provided by the instantiating class itself.
This way the code inside a mixin can be written in an abstract way to promote modularity.

rstrategies supports two types of collections: variable sized and fixed sized lists.
These two abstract types of collections cover the needs of primitive collections for most programming languages.
The reasons behind this design decision are further explained in section \ref{sec:discussion}.

The rstrategies library provides a hierarchy of mixins that can be combined into actual collection classes.
If required, parts of the provided functionality can instead be implemented in the VM specific collection classes.
Any mixin class can also be extended in the VM specific code before being used.
There is one main mixin hierarchy that provides the largest part of the functionality.
At the top of this hierarchy is the class \lstinline{AbstractStrategy}.

\lstinline{AbstractStrategy} defines the interface implemented in its specific subclasses, in addition to a few internal methods.
The following methods are used to interact with the various storage strategy implementations:
\begin{description}
\item[store(index, item)]{Stores the given \lstinline{item} at the given \lstinline{index} in the collection. 
The strategy will check internally if it can store the value of \lstinline{item} in an optimized way.
If this check fails, the strategy will be automatically switched to a more general strategy that is able to hold both the current elements and the new \lstinline{item}.
The switching mechanism is explained in section \ref{sec:switching}.
Optionally, the value of \lstinline{index} is checked and an error is raised if it is out of bounds of the collections size.
The value of \lstinline{index} must be within the bounds of the collection, even if the collection is variable sized.
}
\item[fetch(index)]{Returns the item stored at the given \lstinline{index} in the collection.
Depending on the underlying strategy, it might be necessary to box the value on the fly to represent it on the VM level.
Because of this, the results of two identical calls to \lstinline{fetch} might have different object identities, while still ayways containing equal values.
The distinction between object identity and value equality is important.
The same goes for a call to \lstinline{fetch} after switching the unterlying storage strategy: the results must contain equal values, but there are no guarantees about the object identity.
}
\item[size()]{Returns the number of elements stored in the collection.
In case of a variable sized collection, the result of this can change after an intermediate call to \lstinline{insert} or \lstinline{delete}.}
\item[insert(index, items)]{Inserts the given list of \lstinline{items} at the given \lstinline{index} in the collection.
This increases the size of the collection by the number of inserted items and should therefore only be used in the context of variable sized collections.}
\item[delete(start, end]{Deletes the items at indices \lstinline{start} (inclusive) to \lstinline{end} (exclusive) from the collection.
All indices must be within the bounds of the collection.
This decreases the size of the collection by \lstinline{$end-start$}.}
\end{description}

\subsection{Strategy Optimizations}
\label{sec:solutionOptimizations}

The interface of \lstinline{AbstractStrategy} is designed in a way that any of its subclasses can be used both in a fixed sized or variable sized context.
A collection becomes fixed sized by simply never invoking any of the \lstinline{insert} or \lstinline{delete} methods.
RPythons optimizer will then determine that underlying storage arrays are never changed in size and accordingly optimize allocations and accesses of such collections.

Different subclasses of \lstinline{AbstractStrategy} represent different storage strategies, as described in chapter \ref{cha:storageStrategies}.
Table \ref{tab:strategiesOverwiew} gives an overview over the supported strategies, which are further explained in the following paragraphs.

\begin{table}
\centering

\begin{tabular}[ht]{ l l }
\toprule
EmptyStrategy & Size always zero \\
SingleValueStrategy & Every element is the same value \\
GenericStrategy & Can hold any element \\
WeakGenericStrategy & Holds on weakly to its elements \\
SingleTypeStrategy & Stores unboxed values of a single type \\
TaggingStrategy & Stores unboxed values of a single type plus one special tag value \\
\bottomrule
\end{tabular}

\caption{Summary of the strategies supported by rstrategies.}
\label{tab:strategiesOverwiew}
\end{table}


The first one is \lstinline{EmptyStrategy}, which represents a collection with zero elements.
The size of a collection using \lstinline{EmptyStrategy} is always zero, the \lstinline{fetch} operation always raises an error, and the \lstinline{store} operation always switches the strategy to a different one.
The \lstinline{EmptyStrategy} can be used together with variable sized lists, like Pythons primitive \lstinline{list} type.
In such a context, the \lstinline{EmptyStrategy} would be the starting point for every empty list.
Then, depending on the first inserted element, the storage strategy would be switched appropriately.
When all elements of a list are removed, the strategy can be optimized to the \lstinline{EmptyStrategy} once again, in case it is reused with a different type of element.
% \citeauthor{strategies} have implemented this strategy for the Python interpreter PyPy.

Another useful strategy mixin is \lstinline{SingleValueStrategy}.
This strategy can only store a single value and therefore does not require to allocate an entire storage array.
Again, this strategy can save a lot of memory when used frequently.
The \lstinline{SingleValueStrategy} must allocate memory only for a single value: the size of the list.
As soon as a different value is stored in this strategy, it must be deoptimized to a more general strategy.
This is useful for programming languages with fixed sized primitive collections that are always instantiated filled with a default element.
For example, a Java \lstinline{Array} is initially always filled with \lstinline{null} values, and a Smalltalk object is initially always filled with \lstinline{nil} values.
Here, the \lstinline{SingleValueStrategy} can be used as initial strategy, before being deoptimized to a more general strategy.

\lstinline{GenericStrategy} is the strategy that is able to handle any element.
It corresponds to a simple, non-optimized list and is the fallback strategy when any other strategy has to be deoptimized.
Depending on the context, there can be cases where a \lstinline{GenericStrategy} is again optimized to another strategy.
For example, if the last element is removed from a variable sized list using a \lstinline{GenericStrategy}, it switches the strategy to \lstinline{EmptyStrategy}.
This way, it will again switch to another optimized strategy next time an element is added.
Similar scenarios can be constructed when using a fixed sized list, but in that case additional information must be maintained about the elements.
For example, keeping track of the number of \lstinline{null} values would allow optimizing a generic fixed-sized array to \lstinline{SingleValueStrategy}.
The associated overhead of maintaining this information would most likely surpass the optimization benefits.

\lstinline{WeakGenericStrategy} is a version of \lstinline{GenericStrategy} that holds on weakly to its elements.
This means that an element of a list using a \lstinline{WeakGenericStrategy} can be garbage collected, if the list contains the last reference to this element.
Many programming languages require support for weak collections or objects on VM level.
Since this strategy serves a very special purpose, collections should never switch from or to this strategy.
Instead, a collection should be instantiated with this strategy and keep it throughout its lifetime.

\lstinline{SingleTypeStrategy} is able to store elements of a single type in an optimized storage array, as described in section \ref{cha:storageStrategies}.
Certain conditions must be met in order to use this strategy.
Namely, the identity of the element objects must be determined only through their value, and element objects must be immutable.
These strict conditions are the same as for pooling objects (see section \ref{sec:pooling}) and are usually only met by primitive types like integer or float datatypes.
Therefore, the \lstinline{SingleTypeStrategy} can be used to implement optimized storage for collections of such primitive datatypes.
The VM has to provide routines for \emph{boxing} and \emph{unboxing} the elements of this strategy.
These routines form the bridge between the machine level and VM level representations of the element values.
For example, assume an \lstinline{IntegerStrategy} has been implemented based on \lstinline{SingleTypeStrategy}.
Now a boxed integer object is added to a collection with such a strategy.
This object has to be unboxed in order to store its value into a machine level array of integers.
When fetching a value from this strategy, the according boxed integer object has to be created again.

The \lstinline{TaggingStrategy} is an extension of \lstinline{SingleTypeStrategy}.
In addition to optimized storage of a single element type, it uses one specific tag value in the value range of unboxed elements to represent some special boxed value.
This is useful when fixed sized collections have a default value for their contents.
For example, a fresh Java \lstinline{Array} is filled with the default value \lstinline{null}.
When this array is filled with \lstinline{Integer} instances, there will often be some slots in the array which still have the value \lstinline{null}.
Now if we were using the \lstinline{SingleTypeStrategy}, we would already have to deoptimize to \lstinline{GenericStrategy} in order to represent both \lstinline{null} and \lstinline{Integer} values in the same collection.
With the help of the \lstinline{TaggingStrategy}, we can avoid this early deoptimization by representing the \lstinline{null} slots with special integer values, for example the maximum possible integer value.

There are two drawbacks to this strategy.
First, it adds the overhead of comparing fetched elements with the tag value.
This is the penalty for representing two different types in the same optimized storage.
Second, the \lstinline{TaggingStrategy} can not store the tag value itself.
Therefore, the tag value should be chosen to occur very rarely in real situations.
This can be hard for integer values, since even the maximum possible integer can very well be produced by the user program.
For float values, the IEEE 754 Standard \cite{ieeeFloat} allows for multiple representations of the \lstinline{NaN} value, one of which can be used as tag value.

The described strategies can further be combined by the VM.
For example, a Smalltalk object can consist of two parts: a fixed sized part containing instance variables, and an optional variable sized part, which is Smalltalks notion of arrays.
Even though the two parts form a single object, they can have different requirements regarding storage strategies.
In this case, the VM can use two instances of different strategies to represent a single object.

The VM can even implement strategies that contain objects with more than one data field.
For example, a language might have a primitive type representing a \emph{point}, consisting of two number fields.
A collection of such \emph{point} objects could be optimized by storing all the point numbers in one consecutive array.
This can be implemented using \lstinline{SingleTypeStrategy} and storing instances of RPythons \lstinline{tuple} type into it.
In this example, the VM would have to provide routines to convert between the tuple of raw data, and an actual \emph{point} object.


\subsection{Switching Strategies}
\label{sec:switching}

An important part of rstrategies is the mechanism that allows a collection to switch between different storage strategies.
The generic and naive way to switch from one strategy to another requires two steps:

\begin{enumerate}
\item{Allocate a new storage strategy of an appropriate type and with an appropriate size.}
\item{Invoke \lstinline{fetch} on the old strategy instance and \lstinline{store} on the new strategy instance for every element in the old instance.}
\end{enumerate}

While the first step can not be avoided, the second step can be largely optimized in many cases.
For example, when switching from a \lstinline{SingleValueStrategy} to a \lstinline{GenericStrategy}, every \lstinline{fetch} operation will yield the same result.
Therefore, an optimized switching routine can be implemented, which allocates a \lstinline{GenericStrategy} and fills its storage array with a fixed value, using a routine similar to \lstinline{memset} of the C standard library.

In rstrategies, such optimized switching routines are implemented using a \lstinline{StrategyFactory}.
An instance of \lstinline{StrategyFactory} knows all available storage strategies, manages the switching of strategies and instantiates new strategies.
The \lstinline{StrategyFactory} class is described in more detail in section \ref{sec:StrategyFactory}.
